% Response to Reviewers — MTAP-D-26-00833 (major revision)
% Left-right Comment | Response tables (ASOC-style), mapped 1:1 to EM decision letter + attachments.
\documentclass[11pt,a4paper]{article}
\usepackage[margin=0.85in]{geometry}
\usepackage{times}
\usepackage{booktabs}
\usepackage{longtable}
\usepackage{array}
\usepackage{hyperref}
\usepackage{xcolor}
\usepackage{colortbl}
\definecolor{headergray}{gray}{0.92}

\begin{document}

\begin{center}
{\Large\bfseries Response to Reviewers}\\[0.4em]
Manuscript ID: MTAP-D-26-00833\\
Title: Decomposition Approaches for Cross-Domain Recommendation Systems\\
Journal: Multimedia Tools and Applications
\end{center}

\vspace{0.5em}
\noindent Dear Editor,\\[0.3em]
Thank you for the major-revision invitation and the detailed reviews. Every comment from the decision letter and the attached Reviewer~3 / Reviewer~7 reports is answered below in Comment\,$\mid$\,Response form. Locations in the revised manuscript are given in brackets.\\[0.35em]
Sincerely,\\
Vibhor Kant\\
Corresponding Author

\vspace{0.4em}
\noindent\textbf{Framing.} This manuscript is a \emph{comparative study} of existing decomposition methods (MF, MFB, NMF, DMF, SVD, plus a matched LightGCN-style GNN baseline). It does \emph{not} propose a new recommendation model. Recommended citations were included only where justifiable for the comparative scope.

\newcommand{\resphead}{%
\rowcolor{headergray}
\textbf{Comment} & \textbf{Response} \\
\midrule
}

%% ========================= Reviewer 1 =========================
\section*{Reviewer 1}
{\footnotesize
\begin{longtable}{@{}>{\raggedright\arraybackslash}p{0.38\textwidth}|>{\raggedright\arraybackslash}p{0.56\textwidth}@{}}
\toprule
\resphead
\endfirsthead
\toprule
\resphead
\endhead
\textbf{Major.} Clarify the $1$--$5$ rating range; add a global-mean ($\mu$) baseline so MAE/RMSE near the scale ceiling can be interpreted.
&
Agreed. All ratings are on a $1$--$5$ scale. A global-mean ($\mu$) baseline is added under the same 5-fold protocol so readers can gauge absolute error levels.
\textit{(Sec.~4 Evaluation Metrics; $\mu$ baseline table)} \\
\midrule
\textbf{Major.} Situate MFB against a modern learning-based recommender; cite Meng et al.\ (ASOC 2025) if feasible.
&
A situating subsection and the Meng et al.\ citation are added. A matched LightGCN-style GNN\_CDRS baseline is also evaluated on the same splits.
\textit{(Relation to stronger CDRS; GNN\_CDRS; MAE/RMSE tables)} \\
\midrule
\textbf{Major.} Note that identical hyperparameters / $E{=}20$ likely under-train DMF and NMF; soften ``intrinsically weaker'' claims.
&
Agreed. We now state that the matched $E{=}20$ protocol may under-train slower methods and soften data-hungry / intrinsically-weaker wording.
\textit{(Result discussion; Conclusions)} \\
\midrule
\textbf{Major.} Add at least one ranking metric (NDCG/Recall); note that only user-overlap is evaluated.
&
User-overlap is now stated as the explicit scope; item-/no-overlap are future work. This revision keeps matched MAE/RMSE; ranking metrics are listed as concrete future work.
\textit{(Problem Formulation; Evaluation Metrics; Conclusions)} \\
\midrule
\textbf{Minor.} Revisit DMF attribution (Xue / He GMF, not Camilli \& M\'ezard).
&
Corrected. DMF is attributed to Xue et al.\ / GMF (He et al.); Camilli \& M\'ezard is kept as a separate theoretical line.
\textit{(Literature Review; DMF\_CDRS)} \\
\midrule
\textbf{Minor.} Present the ``MFB below 1\% density'' rule as a hypothesis, not a demonstrated crossover.
&
Agreed. It is now framed as a hypothesis because all six Amazon targets already sit below 1\% density; $ML_1$ is a denser check.
\textit{(Datasets; Result discussion; Conclusions)} \\
\midrule
\textbf{Minor.} Light English edit; anchor GNN future work to Shu et al.\ (IEEE Access, 2024).
&
First-person-singular phrasing and capitalization inconsistencies were cleaned. GNN future work is anchored to Shu et al.\ (2024) with concrete next steps.
\textit{(Literature Review; Conclusions)} \\
\bottomrule
\end{longtable}
}

%% ========================= Reviewer 3 =========================
\section*{Reviewer 3}
\noindent\textit{Responses follow the attached report (I Major / II Minor / IV suggested references).}
{\footnotesize
\begin{longtable}{@{}>{\raggedright\arraybackslash}p{0.38\textwidth}|>{\raggedright\arraybackslash}p{0.56\textwidth}@{}}
\toprule
\resphead
\endfirsthead
\toprule
\resphead
\endhead
\textbf{I-1a.} Abstract: add an explicit state-of-the-art positioning sentence.
&
Added: recent CDRS emphasize complex DL/GNN/adversarial models while foundational transfer behaviour under controlled sparsity remains under-examined.
\textit{(Abstract)} \\
\midrule
\textbf{I-1b.} Abstract: include at least one quantitative performance figure.
&
Added a representative quantitative statement (MFB MAE and relative gain vs.\ matched baselines under the reported protocol).
\textit{(Abstract)} \\
\midrule
\textbf{I-2a.} Introduction: strengthen domain importance with quantitative / platform-scale context.
&
Introduction now includes brief adoption/impact context for recommendation-driven platforms, with supporting references.
\textit{(Section 1)} \\
\midrule
\textbf{I-2b.} Introduction: describe application domains with concrete examples (rating / implicit / social).
&
E-commerce rating, entertainment implicit-feedback, and social-graph examples are briefly illustrated.
\textit{(Section 1)} \\
\midrule
\textbf{I-3a.} Expand related-work table (year, model, datasets, reported metrics where available).
&
Related-work table was expanded within space limits; survey/method rows are separated and scoped to CDRS-relevant work.
\textit{(Section 2 related-work table)} \\
\midrule
\textbf{I-3b.} Add a structured taxonomy / family classification of CDRS techniques.
&
A taxonomy paragraph organizes families (decomposition, deep, adversarial, GNN, pretrain/meta) with representative methods.
\textit{(Section 2)} \\
\midrule
\textbf{I-3c.} Discuss advantages and limitations of each technique family.
&
A family-level pros/cons discussion is added to justify focusing on foundational decompositions.
\textit{(Section 2)} \\
\midrule
\textbf{I-3d.} Cite/discuss negative transfer in CDRS.
&
Negative transfer is discussed and linked to the observed NMF behaviour under non-negativity constraints.
\textit{(Section 2; Result discussion)} \\
\midrule
\textbf{I-4a.} Sharpen the precise problem statement / scope of the comparative answer.
&
Problem Formulation now states measurable failure modes under sparsity, why complex models are not a substitute for this baseline study, and the paper's comparative scope.
\textit{(Problem Formulation / Section 3)} \\
\midrule
\textbf{I-4b.} Provide theoretical time-complexity for each method.
&
Asymptotic complexities for MF/MFB/NMF/SVD/DMF (and GNN) are reported and compared.
\textit{(Method comparison / complexity table)} \\
\midrule
\textbf{I-4c.} Show hyperparameter sensitivity ($d,\lambda,\eta,\gamma$).
&
Sensitivity results for $\gamma$ and $d$ (and related settings) are reported; chosen values are justified around stable regions.
\textit{(Sensitivity analysis)} \\
\midrule
\textbf{I-4d.} Detail DMF architecture (layers, activation, regularization).
&
DMF\_CDRS is specified as a shallow GMF-style interaction (Hadamard + linear) with regularization details for fair matched comparison.
\textit{(DMF\_CDRS; Implementation Details)} \\
\midrule
\textbf{I-4e.} Add formal Algorithm / pseudocode blocks for all methods.
&
Algorithms for MF, MFB, NMF, SVD, DMF, and GNN\_CDRS training/transfer are included.
\textit{(Section 3 Algorithms)} \\
\midrule
\textbf{I-4f.} Justify studying only the user-overlap scenario.
&
User-overlap is justified as the most practical Amazon co-purchase setting; item-/no-overlap are acknowledged as limitations / future work.
\textit{(Problem Formulation; Data pre-processing; Conclusions)} \\
\midrule
\textbf{I-5a.} Compare against at least one stronger / complex CDRS baseline on the same splits.
&
A matched LightGCN-style GNN\_CDRS baseline is added on the same Amazon/$ML_1$ protocol.
\textit{(GNN\_CDRS; MAE/RMSE tables)} \\
\midrule
\textbf{I-5b.} Report wall-clock runtime / scalability evidence.
&
Wall-clock runtimes on a representative pair ($CX_1$) are reported alongside complexity.
\textit{(Runtime table)} \\
\midrule
\textbf{I-5c.} Specify hardware/software stack and random seeds.
&
Colab CPU/High-RAM environment, key library versions, and seed settings are documented.
\textit{(Implementation Details / Experimental setup)} \\
\midrule
\textbf{I-5d.} Add a conceptual comparison of what each method transfers.
&
A transfer-comparison table covers knowledge type, optimization, constraints, complexity, and interpretability.
\textit{(Section 3 conceptual comparison table)} \\
\midrule
\textbf{I-5e.} Justify $E{=}20$ with a convergence study/figure.
&
A convergence figure/study supports the matched iteration budget; under-training risk for slower methods is acknowledged.
\textit{(Convergence figure; Parameter Settings)} \\
\midrule
\textbf{I-6a.} Make future work concrete; ground it in stated limitations.
&
Future work now lists concrete steps: ranking metrics; LightGCN$+$bias; LLM no-overlap; federated transfer; denser/implicit data; Amazon-only / fixed-$d$ limits.
\textit{(Conclusions and Future Works)} \\
\midrule
\textbf{II-1.} Correct Amazon dataset citation (He \& McAuley 2016 / McAuley \& Leskovec).
&
Dataset citation corrected/supplemented to He \& McAuley (2016).
\textit{(Datasets)} \\
\midrule
\textbf{II-2.} Improve readability of Figs.~7--8 (fonts, legends, size).
&
Figure fonts/legends and layout were enlarged for readability.
\textit{(Figs.~7--8)} \\
\midrule
\textbf{IV.} Cite the additional suggested references (incl.\ DisenCDR, Zhang survey, Fern\'andez-Tob\'ias, Shu TKDE 2018, Briels 2024) where justifiable.
&
All suggested references are cited in context (dataset, CDRS survey/methods, identity linkage, and the Briels reproducibility note as requested).
\textit{(Section 2; Datasets; Implementation Details)} \\
\bottomrule
\end{longtable}
}

%% ========================= Reviewer 4 =========================
\section*{Reviewer 4}
{\footnotesize
\begin{longtable}{@{}>{\raggedright\arraybackslash}p{0.38\textwidth}|>{\raggedright\arraybackslash}p{0.56\textwidth}@{}}
\toprule
\resphead
\endfirsthead
\toprule
\resphead
\endhead
1. Novelty missing; motivation unclear; no new concept proposed.
&
Agreed that no new model is proposed. Contribution is reframed as a controlled comparative evaluation of foundational transfer operators under sparsity, with clearer motivation and tempered claims.
\textit{(Abstract; Introduction; Conclusions)} \\
\midrule
2. Only a single dataset family is used.
&
Amazon $CX_1$--$CX_6$ remain the primary sparse family; a denser MovieLens-1M $ML_1$ (Action$\to$Adventure) check is added under the same protocol.
\textit{(Datasets; MovieLens $ML_1$)} \\
\midrule
3. Amazon is not justified as a benchmark.
&
Amazon is justified as a widely used explicit-rating sparse CDRS benchmark (He \& McAuley), suited to stress-testing transfer under extreme sparsity.
\textit{(Datasets)} \\
\midrule
4. Results weak: no statistical analysis; no component-level measures.
&
Paired $t$-tests / CIs are reported. Component-level cold/mid/warm and rating-bin MAE, plus an MFB bias ablation, are added on $CX_1$.
\textit{(Statistical significance; Cold-start / rating-bin; Ablation)} \\
\midrule
5. Consider datasets from various domains to enrich evaluation.
&
$ML_1$ provides a second domain family; broader multi-platform evaluation remains future work and is acknowledged as a limitation.
\textit{(Datasets; Conclusions)} \\
\bottomrule
\end{longtable}
}

%% ========================= Reviewer 5 =========================
\section*{Reviewer 5}
{\footnotesize
\begin{longtable}{@{}>{\raggedright\arraybackslash}p{0.38\textwidth}|>{\raggedright\arraybackslash}p{0.56\textwidth}@{}}
\toprule
\resphead
\endfirsthead
\toprule
\resphead
\endhead
\textbf{Major 1.} Writing quality / grammar / typesetting need a thorough edit.
&
Introduction, Literature Review, and Method sections were edited for grammar, incomplete sentences, capitalization, and equation presentation.
\textit{(Whole manuscript)} \\
\midrule
\textbf{Major 2.} Soften ``first'' / novelty claims; state the contribution more precisely.
&
Contribution language now emphasizes the controlled sparsity sweep and bias-transfer finding, not primacy claims.
\textit{(Abstract; Introduction; Conclusions)} \\
\midrule
\textbf{Major 3.} Fairness: DMF is shallow; transfer mechanisms differ; show sensitivity; only user-overlap studied.
&
DMF is explicitly a shallow GMF-style matched baseline. Transfer differences are analyzed conceptually and via MFB ablation. $\gamma$/$d$ sensitivity is shown. User-overlap scope is justified; other overlaps are future work.
\textit{(DMF\_CDRS; Ablation; Sensitivity; Problem Formulation)} \\
\midrule
\textbf{Major 4.} Complete CX4--CX6 tables; add significance tests; runtime; isolate cold-start users/items.
&
All six Amazon pairs are fully reported. Paired tests, runtime, and cold/mid/warm MAE on $CX_1$ are added.
\textit{(MAE/RMSE tables; Statistical significance; Runtime; Cold-start)} \\
\midrule
\textbf{Major 5.} Literature table incomplete; add at least a short comparison vs stronger modern baselines.
&
Related-work table cleaned/expanded. Matched GNN\_CDRS contextualizes absolute performance against a stronger baseline class.
\textit{(Section 2; GNN\_CDRS)} \\
\midrule
\textbf{Minor.} Notation table incomplete; improve Fig.~7--8 labels; discuss denser-domain generalization; make future work less generic and link bias-transfer to GNN/LLM directions.
&
Notation updated; figure labels improved; denser-domain hypothesis discussed via $ML_1$; future work now gives concrete GNN$+$bias and LLM no-overlap steps.
\textit{(Notation; Figs.~7--8; Datasets; Conclusions)} \\
\bottomrule
\end{longtable}
}

%% ========================= Reviewer 6 =========================
\section*{Reviewer 6}
\noindent No comments were provided by Reviewer~6.

%% ========================= Reviewer 7 =========================
\section*{Reviewer 7}
\noindent\textit{Responses follow the attached revision points.}
{\footnotesize
\begin{longtable}{@{}>{\raggedright\arraybackslash}p{0.38\textwidth}|>{\raggedright\arraybackslash}p{0.56\textwidth}@{}}
\toprule
\resphead
\endfirsthead
\toprule
\resphead
\endhead
1. Clarify novelty vs existing CDRS studies; what is new in this systematic evaluation.
&
Novelty is stated as a controlled sparsity sweep of foundational transfer operators (not a new architecture), clearly differentiated from complex CDRS papers.
\textit{(Abstract; Introduction; Relation to stronger CDRS)} \\
\midrule
2. Expand dataset description (users/items/ratings/density) with a detailed table.
&
A consolidated statistics table for $CX_1$--$CX_6$ and $ML_1$ is included.
\textit{(Dataset statistics table)} \\
\midrule
3. Justify the minimum-30 filter and discuss its effect on sparsity/characteristics.
&
The min-30 filter is justified for stable CV splits; resulting denser-than-raw matrices are acknowledged.
\textit{(Data pre-processing)} \\
\midrule
4. Clarify how 20/40/60/80\% training levels relate to target sparsity / evaluation.
&
TR fractions are defined as supervision levels on the filtered target, distinct from catalogue density; their effect on evaluation is clarified.
\textit{(Experiment Parameter Settings and Protocols)} \\
\midrule
5. Describe hyperparameter selection / sensitivity more clearly.
&
Selection via sensitivity experiments is described; $\gamma$/$d$ results are shown for reproducibility.
\textit{(Parameter Settings; Sensitivity)} \\
\midrule
6. Strengthen statistical analysis of method differences.
&
Paired $t$-tests and 95\% CIs accompany ranking claims.
\textit{(Statistical significance; MAE/RMSE tables)} \\
\midrule
7. Examine why MFB wins; ablate transferred bias terms.
&
An ablation isolates modelling/transferring biases as the main driver under this protocol.
\textit{(Ablation table)} \\
\midrule
8. Explain SVD vs MF/MFB differences (initialization vs continuous transfer).
&
SVD uses truncated warm-start; MF/MFB use continuous $\gamma$-regularization; this distinction is clarified in methods and discussion.
\textit{(SVD\_CDRS; Result discussion; Method comparison)} \\
\midrule
9. Acknowledge Amazon-only / selected-setting limitations and generalization.
&
Limitations now state Amazon primary family, single $ML_1$ check, user-overlap only, and matched hyperparameters.
\textit{(Conclusions and Future Works)} \\
\midrule
10. Language editing; reduce repetition; temper unsupported strong claims.
&
Grammar/readability were improved; repeated statements reduced; strong claims tied directly to reported experiments.
\textit{(Whole manuscript)} \\
\bottomrule
\end{longtable}
}

\vspace{0.6em}
\noindent We hope the revised comparative manuscript and this point-by-point response address all reviewers' concerns.

\end{document}
